\subsection{Synthea Data Ingestion and ETL Pipeline}
\label{ch:synthea_design}

The implementation of any of the algorithms described in this paper would not have been possible without the presence of an adequate dataset. An open-source Synthea model was chosen to serve that purpose due to its ability to provide realistic synthetic data on patients. Of the numerous files generated by the program, several administrative/financial datasets, namely \texttt{claims}, \texttt{claims\_transactions}, \texttt{payers}, \texttt{payer\_transitions}, \texttt{devices}, \texttt{imaging\_studies}, and \texttt{supplies}, were consciously omitted from the analysis process. The exclusion of these specific tables was done based on the fact that they address financial aspects of health care, which makes them irrelevant to our work.

\begin{figure}[H]
\centering
\includegraphics[width=0.8\linewidth]{figs/chapter_5/etl_pipeline.png}
\caption{ETL Pipeline Diagram}
\label{fig:etl_pipeline}
\end{figure}

\clearpage

Normalization is among the key concerns that go into designing the output relational database schema. Each table possesses its own well-defined primary key, and all non-key attributes are fully dependent on the entire key, not some part of it. Many-to-many relations are dealt with properly, employing junction tables, as exemplified by \texttt{patient} and \texttt{provider} connected via \texttt{patient\_provider} table.

Several calculated columns are provided in the database: age of patient is computed based on \texttt{birthdate}, duration of medication and careplan is derived from corresponding start and end date, \texttt{organization.utilization}, \texttt{provider.encounters}, and \texttt{provider.procedures} are counts. They remain in the database on purpose, in order to save the processing power on aggregating and joining data for queries, a common optimization technique for an analytical database in which reading is more frequent than writing.

